\section{Charge-preserving reduction to strict rest data}
\label{ch:rest:section}

Throughout this section the closed forms $\alpha_1,\alpha_2$ are fixed.
For each new metric they define new electromagnetic vector fields by
contraction with its volume form. Thus every change below preserves both
Maxwell constraints and both signed charges. We first apply the
neutral endpoint corollary to establish timelikeness of the original data.

\begin{proposition}[Strict rest reduction]
\label{ch:rest:reduction}
For a charged exterior as in Definition~\ref{ch:pre:charged-exterior}, $E>|P|$ and
$m=(E^2-|P|^2)^{1/2}>0$. There are smooth data $(g_j,K_j)$ on the same
exterior, with the same flux forms, satisfying the following properties.
\begin{enumerate}
\item The physical matter condition is strict everywhere, the original
boundary $S$ has $\theta_{+,j}<0$, and the constraint densities remain
integrable. For one fixed $0<\delta<1$ the matter margin satisfies
\[
 M_{0,j}=8\pi(\mu_{m,j}-|J_{m,j}|_{g_j})
       \ge c_j r_y^{-3-\delta}\quad\hbox{far out},\qquad c_j>0.
\]
\item In a rest chart $y$, $K_j$ is compactly supported and
\[
 g_j-\delta_{\mathrm{eucl}}=O_k(r_y^{-1}),\qquad
 R_{g_j}=O_k(r_y^{-3-\delta})\quad\hbox{for every }k.
\]
The fields have the required $O_1(r_y^{-2})$ decay. In particular the
ADM momentum is zero and $\Ric_{g_j}=O(r_y^{-3})$.
\item Smooth convergence holds on every compact set, including at $S$,
and
\begin{equation}
 \label{ch:rest:limits}
 E_j\longrightarrow m,\qquad
 a_{g_j}(S)\longrightarrow a_g(S),\qquad
 (Q_{E,j},Q_{B,j})=(Q_E,Q_B).
\end{equation}
Moreover $g_j\ge(1-\varepsilon_j)g$ globally for some
$\varepsilon_j\downarrow0$.
\end{enumerate}
\end{proposition}

The proof has four steps: establish timelikeness from the neutral
endpoint corollary; replace the distant end by a vacuum reference with
matching ADM vector; bend and repair while retaining the flux forms;
and make the inequalities strict by a separate elliptic construction.
All quantitative estimates for the replacement use exactly the original
$O_2$ metric and $O_1$ tensor assumptions.

\subsection{Timelikeness}

\begin{lemma}
\label{ch:rest:timelike}
The original ADM vector satisfies $E>|P|$.
\end{lemma}
\begin{proof}
The physical matter condition and Equation~\eqref{ch:pre:neutral-margin}
give
\[
 \mu-|J|_g\ge\mu_m-|J_m|_g\ge0.
\]
These are the actual total constraint densities of $(g,K)$.
Definition~\ref{ch:pre:charged-exterior} supplies smoothness, orientation,
one-endedness, completeness with the nonempty compact boundary included,
the required metric and tensor falloff, integrable total densities,
finite ADM limits, and weak future trapping with the normal toward the
end. The total constraint and ADM conventions are exactly those of
Equations~\eqref{n3:intro:constraints}, \eqref{n3:intro:energy} and
\eqref{n3:intro:momentum}. All hypotheses of
Corollary~\ref{bridge:bridge:timelike} are therefore satisfied.
Applying it to the unchanged data $(N,g,K)$ yields $E>|P|$
for their original ADM vector.
\end{proof}

\subsection{A matching vacuum reference and annular replacement}

\begin{lemma}[Boosted vacuum reference]
\label{ch:rest:reference}
There is an auxiliary Schwarzschild end of mass $m$ whose induced data
$(g_*,K_*)$, in asymptotically Euclidean coordinates $x$, have ADM vector
$(E,P)$ and satisfy
$g_*-\delta_{\mathrm{eucl}}=O_k(r_x^{-1})$,
$K_*=O_k(r_x^{-2})$ for every $k$.
\end{lemma}
\begin{proof}
Lemma~\ref{ch:rest:timelike} gives $E>|P|$. The constraint and ADM
conventions in Equations~\eqref{ch:pre:constraints},
\eqref{ch:pre:adm-energy}, and~\eqref{ch:pre:adm-momentum} agree with
those in Part~I. Lemma~\ref{n3:reduction:boosted-reference} therefore
gives the required reference end, including the stated differentiated
decay and the future second-form convention.

For the later bend, retain its static isotropic coordinates:
\begin{equation}
 \label{ch:rest:schwarzschild}
 \mathbf g=-\left(\frac{1-m/(2r_y)}{1+m/(2r_y)}\right)^2dT^2
       +\left(1+\frac m{2r_y}\right)^4|dy|^2.
\end{equation}
The reference plane is $T=v\cdot y$, with $|v|<1$ chosen to give
$(E,P)$. Its asymptotically Euclidean coordinates are $y=Ax$, where
$A=(I-vv^\top)^{-1/2}$, since $A^\top(I-vv^\top)A=I$.
\end{proof}

For the remaining construction, fix $1/2<\beta<\min(q,1)$.

\begin{lemma}[Annular replacement]
\label{ch:rest:annular}
For sufficiently large $R$ there are smooth data equal to $(g,K)$ on
$r_x\le R$ and to $(g_*,K_*)$ on $r_x\ge4R$, with possible neutral
DEC deficit bounded by $\eta_R R^{-3}$ in between, where
$\eta_R\to0$. The estimates use no derivatives beyond the assumed
two metric derivatives and one tensor derivative.
\end{lemma}
\begin{proof}
The original data satisfy neutral DEC by
Equation~\eqref{ch:pre:neutral-margin}. Definition~\ref{ch:pre:charged-exterior}
gives the smooth one-ended exterior, completeness with its compact
boundary included, decay $g-\delta=O_2(r_x^{-q})$ and
$K=O_1(r_x^{-1-q})$ with $q>1/2$, integrable total constraint
densities, and finite ADM limits. The total constraint and ADM
normalizations agree with those used in
Proposition~\ref{n3:reduction:annular-replacement}.
The reference in Lemma~\ref{ch:rest:reference} has matching charges
in the common $x$ coordinates, and $E>|P|$ by
Lemma~\ref{ch:rest:timelike}. That proposition applies and gives
\[
 16\pi(\widetilde\mu_R-
       |\widetilde J_R|_{\widetilde g_R})
 \ge -\eta_RR^{-3},\qquad \eta_R\longrightarrow0,
\]
with the stated exact agreement on both sides of the annulus and
with quantitative estimates using only the original two metric
derivatives and one tensor derivative. Absorbing the fixed
normalization into $\eta_R$ gives the assertion.
The scaled construction in Equation~\eqref{n3:reduction:scaled-smallness}
and the compact corrections in that proposition's proof also give
\[
 \|\widetilde g_R(R\,\cdot)-I\|_{C^2(\{1<|z|<4\})}
 +\|R\widetilde K_R(R\,\cdot)\|_{C^1(\{1<|z|<4\})}
 \le C R^{-\beta}.
\]
The two closed flux forms are retained throughout; their contribution
to the physical matter deficit is estimated next.
\end{proof}

\subsection{Bending and repairing the physical matter condition}

Beyond the annular replacement, use the exact Schwarzschild part and
replace the plane by the graph
\begin{equation}
 \label{ch:rest:bend}
 T_R(y)=(v\cdot y)\psi\left(\frac{\log(r_y/R_b)}{L_0}\right).
\end{equation}
Here $\psi$ is smooth, is one before zero and zero after one, and takes
values in $[0,1]$. The fixed multiple $R_b/R$ is large enough that the
bend begins beyond $r_x=4R$. First fix $L_0$ so large that
$|v|(1+\|\psi'\|_\infty/L_0)<1$. Then the graph is uniformly
spacelike: in Equation~\eqref{ch:rest:schwarzschild} the spatial factor
is greater than one and the lapse is less than one. Its constraints
vanish by Gauss--Codazzi with the prescribed future normal.

Write $(g_R,K_R)$ for the unscaled replacement and bend. Constants in
the rest of this subsection may depend on the fixed original data and
$L_0$, but not on $R$. The graph formula and its first three derivatives
give uniform metric comparability, $|\partial g_R|+|K_R|\le C/r_y$ in
the transition, and
\begin{equation}
 \label{ch:rest:radial-geometry}
 c\le|dr_y|_{g_R}\le C,\qquad
 |\Delta_{g_R}r_y|\le C/r_y.
\end{equation}
The same bounds hold on the gluing annulus, since $y=Ax$ is a fixed
linear map. Outside a fixed multiple of $R$, the data are horizontal
Schwarzschild and $K_R=0$.

Keep the original closed forms in these coordinates. Their norms are
$O(r_y^{-2})$, so their contribution to the negative part of
$\inf_{|Y|\le1}D_{g_R}(Y)$ is $O(r_y^{-4})$. Choose fixed
$0<a<\xi_0<b$ so that the data are original for $r_y/R\le\xi_0$
and horizontal Schwarzschild for $r_y/R\ge b$.
Lemma~\ref{ch:rest:annular} and vacuum of the reference and bend imply
\begin{equation}
 \label{ch:rest:charged-deficit}
 \bigl[-\inf_{|Y|_{g_R}\le1}D_{g_R}(Y)\bigr]_+
 \le
 \begin{cases}
 0,& r_y/R\le\xi_0,\\
 \eta_R R^{-3},&\xi_0\le r_y/R\le b,\\
 Cr_y^{-4},&r_y/R\ge b,
 \end{cases}
 \qquad\eta_R\longrightarrow0.
\end{equation}
Thus this intermediate data set is not assumed to satisfy charged DEC.

\begin{lemma}[Charged repair with vanishing energy cost]
\label{ch:rest:charged-repair}
There are smooth $u_R\ge1$, constant on the unchanged core, for which
$(u_R^4g_R,u_R^2K_R)$ with the fixed forms satisfy the physical matter
condition. They retain weak trapping, have zero ADM momentum, and their
energies tend to $m$. Moreover $\|u_R-1\|_\infty=o(R^{-1})$, and
the final metric has all-order $O(r_y^{-1})$ asymptotics and
scalar curvature $R_{u_R^4g_R}=O_k(r_y^{-4})$.
\end{lemma}
\begin{proof}
Set $\xi=r_y/R$ and seek
\[
 u_R=1+\frac{e_R}{R}F_R(\xi),\qquad
 F_R(\xi)=\int_\xi^\infty z_R(t)\,dt,
 \qquad z_R\ge0,
\]
where $e_R\to0$ will be chosen below. On $[a,b+1]$ initially solve
$z_R'=C_0z_R+\omega$, $z_R(a)=0$, where $\omega\ge0$ is smooth,
zero near $a$, and at least one on $[\xi_0,b+1]$; extend $z_R=0$
below $a$. Equation~\eqref{ch:rest:radial-geometry} yields
\begin{align}
 -\Delta_{g_R}u_R-|K_R|_{g_R}|du_R|_{g_R}
 &=\frac{e_R}{R^3}
  \{z_R'|dr_y|^2+Rz_R\Delta r_y-Rz_R|K_R||dr_y|\}\notag\\
 &\ge\frac{e_R}{R^3}(c z_R'-Cz_R).
 \label{ch:rest:repair-inequality}
\end{align}
Choose $C_0$ large. This is nonnegative on $[a,b]$ and at least
$c'e_R R^{-3}$ on $[\xi_0,b]$.

On the horizontal region put
$A_R(t)=(t+m/(2R))^2z_R(t)$. Its derivative from the preceding
prescription is positive near $b$. Continue that derivative smoothly
by a convex blend with $t^{-2}$ over $[b,b+1]$, and integrate from
$A_R(b)$. Thus $A_R'\ge c''t^{-2}$ for $t\ge b$ and
$A_R'=t^{-2}$ eventually. Defining
$z_R=A_R(t)/(t+m/(2R))^2$ there makes every join smooth. For fixed
$L_0$, $A_R^\infty=\lim_{t\to\infty}A_R(t)$ and $F_R$ are bounded
uniformly in $R$. The exact radial Laplacian on the horizontal slice is
\begin{equation}
 \label{ch:rest:tail-repair}
 -\Delta_{g_R}u_R
 =\frac{e_R}{R^3}
   \left(1+\frac m{2R\xi}\right)^{-6}\xi^{-2}A_R'(\xi)
 \ge c'''e_R R^{-3}\xi^{-4}.
\end{equation}
Choose, for example, $e_R=C_1(\eta_R+R^{-1})$, with a sufficiently
large fixed $C_1$. For large $R$, $1\le u_R\le2$.
Equations~\eqref{ch:rest:repair-inequality}--\eqref{ch:rest:tail-repair}
then dominate $u_R/4$ times the deficit in
Equation~\eqref{ch:rest:charged-deficit}.
Equation~\eqref{ch:pre:conformal-estimate} proves the physical matter condition.
The factor is constant near $S$, so
Equation~\eqref{ch:pre:conformal-boundary} preserves its weak expansion.

At infinity, for each fixed $R$,
\[
 u_R=1+\frac{e_RA_R^\infty}{r_y}+O_k(r_y^{-2})
 \quad\hbox{for every }k.
\]
The ADM energy is therefore $m+2e_RA_R^\infty\to m$ and the momentum
vanishes because $K_R$ has compact support. The differentiated radial
formula gives scalar curvature $O_k(r_y^{-4})$. This also verifies
integrability of the repaired constraints. Completeness follows from
the comparability of the intermediate metric and $u_R\ge1$.
\end{proof}

\subsection{Strictification by an exterior elliptic problem}

\begin{lemma}[A strict conformal direction]
\label{ch:rest:strict-direction}
Fix one repaired exterior $(g_0,K_0)$ from
Lemma~\ref{ch:rest:charged-repair}, and fix $0<\delta<1$. There exist
smooth positive $w_0,\phi_0$ such that
\begin{equation}
 \label{ch:rest:strict-direction-equations}
 \partial_\nu\phi_0=-1\text{ on }S,\qquad
 -\Delta_{g_0}\phi_0-|K_0|_{g_0}|d\phi_0|_{g_0}\ge w_0,
\end{equation}
where $w_0=r_y^{-3-\delta}$ and
$-\Delta_{g_0}\phi_0=w_0$ far out, and
$\phi_0=O_k(r_y^{-1})$ for every $k$.
Its gradient has finite flux at infinity. All constants in this lemma
may depend on the fixed repaired exterior.
\end{lemma}
\begin{proof}
Choose a smooth nonnegative compactly supported majorant
$b_0\ge|K_0|_{g_0}$ and a smooth positive extension $w_0$ of the
specified tail. We solve the regularized equation
\begin{equation}
 \label{ch:rest:semilinear}
 -\Delta_{g_0}\phi=b_0\sqrt{1+|d\phi|_{g_0}^2}+w_0,
 \qquad\partial_\nu\phi=-1,\qquad\phi\longrightarrow0.
\end{equation}
First truncate at $r_y=T$ and impose outer Dirichlet value zero.
Multiply $b_0$ by $t\in[0,1]$ for continuation. At $t=0$ the mixed
Laplacian is invertible, by the Dirichlet Poincar\'e inequality and
elliptic regularity. Every linearization is a uniformly elliptic
Laplacian with drift of norm at most $b_0$, with homogeneous Neumann
data on $S$ and Dirichlet data on the outer sphere. These faces are
disjoint. The maximum and Hopf principles give zero kernel; as a
compact lower-order perturbation of the mixed Laplacian, it is
Fredholm of index zero and is invertible.

We supply the estimates needed for closedness. On this fixed truncation,
the $W^{2,p}$ estimate and gradient interpolation give, uniformly in $t$,
\[
 \|\phi\|_{W^{2,p}}
 \le C(1+\|\phi\|_{L^p}+\|d\phi\|_{L^p})
 \le\tfrac12\|\phi\|_{W^{2,p}}+C'(1+\|\phi\|_\infty).
\]
If the supremum norms were unbounded, divide by them and take $p>3$.
A subsequence converges in $C^1$ to a nonzero solution $v$ with
homogeneous mixed data and
$-\Delta_{g_0}v=t_*b_0|dv|$. Indeed the normalized right side is
$t b_0\sqrt{M^{-2}+|dv|^2}+w_0/M$ when $M=\|\phi\|_\infty$.
The limit is a homogeneous bounded-drift equation, taking drift
$t_*b_0\nabla v/|dv|$ off the critical set and zero on it. The strong
and Hopf maximum principles exclude such a nonzero $v$.
The resulting supremum and $W^{2,p}$ bounds, followed by Schauder
estimates for the smooth regularized equation, close continuation.
Thus every sufficiently large truncation has a smooth solution.

All solutions are nonnegative: a negative minimum cannot be interior
since the source is positive, or on $S$ since the outward-domain
derivative there is $-\partial_\nu\phi=1$, contradicting the Hopf
sign at a minimum. Fix a core containing $\supp b_0$. On its exterior
$f=r_y^{-1}-r_y^{-1-\delta}$ is positive and
\[
 -\Delta_{g_0}f
 =\delta(1+\delta)r_y^{-3-\delta}+O(r_y^{-4})
 \ge c w_0
\]
after enlarging the fixed core. Comparison with a multiple of $f$
gives, independently of $T$,
\begin{equation}
 \label{ch:rest:core-barrier}
 0\le\phi\le C(1+M_T)r_y^{-1}\quad\hbox{outside the core},
 \qquad M_T=\sup_{\mathrm{core}}\phi.
\end{equation}
If $M_T\to\infty$ along an exhaustion, divide by $M_T$. Local
$W^{2,p}$ estimates, including the fixed inner Neumann boundary,
and Equation~\eqref{ch:rest:core-barrier} give a limit $v\ge0$ with
maximum one on the core, homogeneous inner Neumann condition,
$-\Delta v=b_0|dv|$, and $v\to0$ at infinity. Its positive global
maximum is attained, contradicting the strong or Hopf principle.
Hence the core bounds are uniform in $T$. Diagonal compactness and
bootstrap give a smooth global solution of
Equation~\eqref{ch:rest:semilinear}, positive also on $S$ by the same
minimum argument.

On the end its equation is the linear Poisson equation
$-\Delta_{g_0}\phi_0=w_0$. On a scaled annulus the function
$\rho\phi_0(\rho\,\cdot)$ is uniformly bounded; the scaled metric
has all-order bounds and the scaled source is $O(\rho^{-\delta})$
with all its derivatives. Interior Poisson estimates of each order
therefore give $\phi_0=O_k(r_y^{-1})$. The full right side of
Equation~\eqref{ch:rest:semilinear} is integrable. The divergence theorem
then gives a finite metric gradient flux, and its difference from the
Euclidean flux is $O(r_y^{-1})$. This proves every assertion.
\end{proof}

For $u=1+s\phi_0$, $s>0$, use $(u^4g_0,u^2K_0)$ and the fixed
forms. Equations~\eqref{ch:pre:conformal-estimate} and
\eqref{ch:rest:strict-direction-equations} imply
\begin{equation}
 \label{ch:rest:strict-margin}
 M_{0,s}\ge4s u^{-5}w_0>0,
 \qquad
 \theta_{+,s}=u^{-2}(\theta_{+,0}-4s/u)<0\quad\hbox{on }S.
\end{equation}
Compact support of $K_0$ persists. The conformal scalar formula gives
the asserted all-order $r_y^{-3-\delta}$ scalar decay, and thus
integrability. The ADM energy change is
\begin{equation}
 \label{ch:rest:strict-energy}
 E_s-E_0=-\frac{s}{2\pi}
   \lim_{r\to\infty}\int_{r_y=r}\partial_{\nu_r}\phi_0\,dA_{g_0}
   =O(s).
\end{equation}
The metric and Euclidean fluxes agree in this formula by the estimates
just proved. For the fixed repaired datum the finite coefficient and
any finite collection of compact smooth norms can be made arbitrarily
small by choosing $s$ small; no bound uniform in the bending radius is
needed.

\subsection{Global area control and the diagonal limit}

\begin{proof}[Completion of the proof of Proposition~\ref{ch:rest:reduction}]
It remains to check area and organize the parameters. On the vacuum
plane and bend, write $dT_R=\psi a+b$, where $a=v\cdot dy$ and
$|b|\le C L_0^{-1}|dy|$. Since $0\le\psi\le1$,
$dT_R^2\le a^2+C L_0^{-1}|dy|^2$ as quadratic forms. The spatial
factor and lapse in Equation~\eqref{ch:rest:schwarzschild} consequently
give
\begin{equation}
 \label{ch:rest:global-lower-comparison}
 g_R\ge |dy|^2-dT_R^2
 \ge|dy|^2-(v\cdot dy)^2-\frac C{L_0}|dy|^2
 \ge\left(1-\frac{C'}{L_0}\right)|dx|^2.
\end{equation}
The constants here depend on the fixed boost, not on $R$ or large
$L_0$. In the gluing annulus $g_R=I+O(R^{-\beta})$, whereas the
original metric is uniformly $I+o(1)$ throughout the receding end.
On the unchanged core the metrics agree. Both conformal changes
increase the metric. Thus first choosing $L_0\to\infty$ and then
$R\to\infty$ sufficiently fast gives the global comparison
$g_j\ge(1-\varepsilon_j)g$ with $\varepsilon_j\to0$.

For each fixed $L_0$ the energy of the repaired data tends to $m$ by
Lemma~\ref{ch:rest:charged-repair}. Select $R$ to make this energy error,
the constant core conformal change, and the preceding comparison error
arbitrarily small. Finally choose the strictification parameter $s_j$
to make Equation~\eqref{ch:rest:strict-energy} tend to zero and to control
the first $j$ smooth norms on the first $j$ sets of a compact exhaustion.
This gives local smooth convergence and all the analytic assertions of
the proposition. Each metric is complete up to $S$, by the global
positive comparison and smoothness. The fixed linear relation $y=Ax$
and the all-order metric asymptotics preserve the $O_1(r_y^{-2})$
field bounds, since the forms themselves are fixed.

For every smooth filled enclosing cut $\Gamma$, the quadratic-form
comparison gives
\[
 \Area_{g_j}(\Gamma)\ge(1-\varepsilon_j)\Area_g(\Gamma).
\]
Hence $\liminf_j a_{g_j}(S)\ge a_g(S)$. Conversely fix any such cut
with area within $\zeta>0$ of $a_g(S)$. Local convergence on this
compact cut gives
$\limsup_j a_{g_j}(S)\le\Area_g(\Gamma)\le a_g(S)+\zeta$;
let $\zeta\downarrow0$. This proves the area limit directly for the
smooth-cut infimum, including obstacle-coincident components.
Finally each enclosing cut has the same flux of each fixed closed
form by Stokes' Theorem, so both charges are preserved exactly.
\end{proof}

The proposition constructs charged initial data on the original exterior
and recovers its invariant mass;
it does not assert that the original data admit a spacetime embedding
or a global Lorentz change of asymptotic slice.
